% !TEX root = ../../main.tex

\section{服务端技术}

\subsection{Spring Boot}

Spring Boot 是基于 Spring 生态的 Java Web 应用快速开发框架，能够通过自动配置、约定式工程结构和丰富的 starter 依赖简化后端服务开发\cite{spring-boot-docs}。
在 HearBridge 系统中，Spring Boot 服务端是连接移动端、后台管理端、数据库、对象存储和 Python 识别服务的核心业务层，负责提供统一的 REST 接口和业务编排能力。

本系统服务端采用 Spring Boot 技术栈实现，主要承担用户认证、手语资源管理、样本管理、模型版本管理、文件上传、识别服务调用和语义修正调用等功能。
相比将所有逻辑直接放在客户端或识别服务中，Spring Boot 服务端作为中间业务层能够更好地统一数据结构、权限校验、错误处理和跨端接口调用。

\subsection{MyBatis}

MyBatis 是一种常用的 Java 持久层框架，支持开发者编写自定义 SQL、存储过程和高级映射关系，并通过 XML 或注解方式将数据库记录映射为 Java 对象\cite{mybatis-docs}。
MyBatis-Spring-Boot-Starter 进一步简化了 MyBatis 与 Spring Boot 项目的集成，使开发者能够更方便地在 Spring Boot 工程中使用 Mapper、SQL 映射和数据访问能力\cite{mybatis-spring-boot-starter}。

在 HearBridge 服务端中，MyBatis 主要用于完成用户信息、手势分类、手势资源、样本数据、模型版本、训练记录等结构化数据的读写操作。
与全自动 ORM 框架相比，MyBatis 允许开发者直接控制 SQL，适合本系统中较多分页查询、条件筛选、状态更新和关联查询场景。

\subsection{MySQL}

MySQL 是应用广泛的关系型数据库管理系统，适合存储结构化业务数据，并支持事务、索引、约束和 SQL 查询等能力\cite{mysql-84-docs}。
HearBridge 系统中的用户信息、手势分类、手势资源、样本元数据、模型版本和训练记录等均属于结构化数据，适合使用关系型数据库进行统一管理。

在本系统中，MySQL 作为核心业务数据库，主要负责保存系统的长期业务数据。
服务端通过 MyBatis 访问 MySQL，使系统核心数据具备较好的持久性和可查询性，也便于后续进行统计分析、训练记录追踪和模型版本回溯。

\subsection{Redis}

Redis 是一种高性能键值型内存数据库，常用于缓存、会话管理、分布式锁和临时数据存储等场景\cite{redis-docs}。
在 HearBridge 系统中，Redis 主要用于存储用户登录态和后台管理员登录态。
用户登录后，服务端生成 token，并将 token 与用户身份信息写入 Redis；
后续移动端或后台管理端请求接口时，通过请求头携带 token，服务端再根据 Redis 中的登录态完成身份校验。

相比将登录状态直接保存在数据库中，Redis 具有读写速度快、过期时间设置方便等优点，更适合保存 token、验证码和临时会话数据。
这样的设计可以降低数据库访问压力，也便于在用户退出登录或登录过期时快速删除对应状态。

\subsection{MinIO}

MinIO 是一种高性能、兼容 S3 协议的对象存储服务，适合存储图片、视频、模型文件和其他非结构化资源\cite{minio-docs}。
HearBridge 系统涉及较多文件资源，包括用户头像、手语资源封面、数字人资源文件、训练样本视频和模型产物文件等。
若将这些文件直接存储在服务端本地目录或数据库中，不利于资源统一管理和后续扩展；
采用对象存储服务能够更清晰地区分结构化业务数据和非结构化文件资源。

在本系统中，MinIO 主要承担文件资源存储和访问支撑作用。
服务端通过文件上传接口接收客户端或管理端上传的图片、视频和资源文件，将文件保存到对象存储中，并将文件对象地址或资源标识写入数据库。
这样的设计降低了文件资源管理复杂度，也为后续迁移到其他对象存储服务提供了可能。

\subsection{OpenAPI 与 Swagger UI}

在前后端分离系统中，接口文档对于移动端、后台管理端和服务端联调具有重要作用。
OpenAPI Specification 是一种面向 HTTP API 的标准化接口描述规范，可用于描述接口路径、请求参数、响应结构和认证方式等内容\cite{openapi-320}。
Swagger UI 则能够根据 OpenAPI 描述生成可交互的接口文档页面，便于开发人员查看和调试接口\cite{swagger-ui}。

在 HearBridge 后端开发过程中，接口文档和调试工具主要用于辅助移动端、后台管理端与服务端之间的联调。
后端项目维护了接口说明文档，并引入 SpringDoc OpenAPI 相关依赖，为接口查看和调试提供辅助支持。
开发过程中可结合接口文档与调试工具验证请求路径、请求参数、请求头和响应结构，从而降低跨端联调成本。
